\chapter{Literature Review}
\section{Review scope and selection}
The review covers research that bears on four questions raised by the project: how the KV cache is structured during generation, how behaviour can be localised to and changed through specific model components, how language-model agents influence one another, and how social bias and factual reliability are measured. Papers were located through the ACL Anthology, the proceedings of NeurIPS, ICLR and ICML, arXiv and Google Scholar, using the keywords \emph{KV cache}, \emph{attention sink}, \emph{KV cache compression}, \emph{heavy hitter}, \emph{causal mediation analysis}, \emph{model editing}, \emph{inference-time intervention}, \emph{function vectors}, \emph{multi-agent language models}, \emph{social bias benchmark} and \emph{truthfulness}. A paper was selected when it was a primary, peer-reviewed publication from 2020 to 2024, when it proposed a method, dataset or analysis directly reused or contrasted in this project, and when its claims could be checked against the published text. Surveys were admitted only where they define the evaluation and mitigation taxonomy used in the project. Eight of the twelve selected works were already cited in the project's design document; the remaining four were added during implementation because they bear directly on the persistence metric, the bias-localisation method and the intervention design.

The selected works fall into four groups. The first group (Sections~\ref{lit:sinks}--\ref{lit:tova}) studies the structure of the KV cache for efficiency. The second group (Sections~\ref{lit:fv}--\ref{lit:iti}) localises behaviour to model components and intervenes on it. The third group (Section~\ref{lit:debate}) studies interaction between language-model agents. The fourth group (Sections~\ref{lit:survey}--\ref{lit:tqa}) defines how bias and truthfulness are measured. Each paper is described in a connected paragraph covering its objective, approach, study context, evaluation, findings, limitations and relevance to this project; numerical values are reported only where they appear in the paper, and details that the paper does not state are marked as not reported.

\section{Descriptive review of literature works}

\subsection{Attention sinks and streaming generation}\label{lit:sinks}
Xiao, Tian, Chen, Han and Lewis (2024) in ``Efficient streaming language models with attention sinks''~\cite{xiao2024sinks} investigate how to deploy language models in streaming applications such as long multi-round dialogue, where the KV cache grows without bound and sliding-window attention collapses once the earliest tokens are evicted. The work identifies the \emph{attention sink} phenomenon, in which a few initial tokens receive large attention scores regardless of their meaning, and proposes StreamingLLM, which keeps the cache entries of these initial tokens together with a rolling window of recent tokens. The study uses the Llama-2, MPT, Falcon and Pythia model families without fine-tuning and additionally shows that pre-training with a dedicated placeholder sink token improves streaming. The authors evaluate language-modelling perplexity on long concatenated texts and streaming question answering, and report stable modelling on sequences of up to four million tokens with a speed-up of up to 22.2 times over sliding-window recomputation. The main limitation is that sinks are studied only as a device for stable and efficient decoding; the method does not extend the usable context, and the semantic content of sink entries and their effect on behaviour are not analysed. This work informs the present project in two ways: it is the reason the attractor analysis holds out the first positions as a separate sink class, and it leaves open whether persistent entries other than sinks influence what the model says.

\subsection{Heavy-hitter oracle for KV cache eviction}\label{lit:h2o}
Zhang, Sheng, Zhou, Chen, Zheng, Cai, Song, Tian, R\'e, Barrett, Wang and Chen (2023) in ``H$_2$O: Heavy-hitter oracle for efficient generative inference of large language models''~\cite{zhang2023h2o} address the memory cost of the KV cache during generation. They observe that a small portion of tokens, the \emph{heavy hitters}, contributes most of the attention mass, that these tokens emerge naturally and correlate with frequent co-occurrence in the text, and that removing them sharply degrades quality. The work formulates cache eviction as a dynamic submodular problem and proposes a greedy policy that keeps recent tokens together with heavy hitters. The study uses OPT, LLaMA and GPT-NeoX models on tasks from the HELM and lm-eval-harness suites, including question answering and summarisation. The authors assess accuracy and generation quality against the cache budget, together with throughput and latency, and report that retaining 20\% of the cache as heavy hitters improves throughput by up to 29 times over DeepSpeed Zero-Inference and Hugging Face Accelerate and by up to 3 times over FlexGen on OPT-6.7B and OPT-30B. The main limitation for the present purpose is that importance is defined purely by accumulated attention and is used only to decide what to evict; no link is made between heavy hitters and behaviours such as bias. The work informs this project through the reuse-score metric, which is a normalised form of accumulated attention, and through the gap it leaves between cache importance and behavioural effect.

\subsection{Persistence of importance in the KV cache}\label{lit:scissor}
Liu, Desai, Liao, Wang, Xie, Xu, Kyrillidis and Shrivastava (2023) in ``Scissorhands: Exploiting the persistence of importance hypothesis for LLM KV cache compression at test time''~\cite{liu2023scissorhands} investigate how to fit the KV cache into a fixed memory budget during inference. The work proposes the \emph{persistence of importance} hypothesis, which states that tokens that receive high attention at one step tend to remain important at later steps, and builds a compression method that keeps such pivotal tokens and drops the others with high probability. The study is conducted on the OPT model family on language-modelling and few-shot downstream tasks. The authors evaluate model quality as a function of the compression ratio and report that the KV cache memory can be reduced by up to five times without loss of quality, and by up to twenty times when combined with 4-bit quantisation. The main limitations are that persistence is examined only as a compression criterion within a single sequence, and that the hypothesis is not tested across turns written by different agents or related to the content held by persistent tokens. The work informs the present project directly: the attractor criterion uses a persistence metric $P(t)$, the fraction of later query and head pairs for which a position stays in the top-$k$ attention, and asks whether persistent entries outside the sink class carry behaviourally relevant information.

\subsection{Transformers as multi-state recurrent networks}\label{lit:tova}
Oren, Hassid, Yarden, Adi and Schwartz (2024) in ``Transformers are multi-state RNNs''~\cite{oren2024msrnn} investigate how decoder-only transformers relate to recurrent neural networks. The work shows that such a transformer can be viewed as a multi-state recurrent network with an unbounded state, the KV cache, and as a finite multi-state network when the cache size is limited. It proposes TOVA (Token Omission Via Attention), a training-free policy that, when the cache is full, drops the entry receiving the least attention. The study evaluates open models including LLaMA-2 and Mistral on long-range tasks such as language modelling over long documents and long-text understanding. The authors assess quality against the cache budget and report that TOVA outperforms other compression policies and often approaches full-cache quality with as little as one-eighth of the original cache. The main limitations are that the analysis is conceptual and efficiency-oriented: it frames the cache as the state of the model but does not track how that state changes semantically during a conversation or how it affects behaviour. The work informs the present project by providing the conceptual justification for treating the KV cache as the runtime state of the model, which is the premise of the title \emph{runtime computational memory}.

\subsection{Function vectors in large language models}\label{lit:fv}
Todd, Li, Sharma, Mueller, Wallace and Bau (2024) in ``Function vectors in large language models''~\cite{todd2024fv} investigate how a language model represents the task demonstrated by in-context examples, such as producing antonyms or translating words. Using causal mediation analysis over a diverse set of in-context learning tasks, the work finds that a small number of attention heads transport a compact representation of the demonstrated task, which the authors call a function vector. The study uses autoregressive models including GPT-J and Llama-2 models of up to 70 billion parameters. The authors evaluate the causal effect of inserting the function vector into the hidden states of a new forward pass, and report that function vectors trigger execution of the task in zero-shot settings and in natural text that does not resemble the original prompts, and that they show some compositional properties. The main limitation is that the analysis is restricted to a single forward pass over one prompt; the representation is not followed through the evolving cache of a multi-turn conversation. The work informs the present project by showing that information obtained from context can be concentrated in a few attention heads, which motivates the search for bias reservoirs, and it leaves open whether such concentrated information persists in cached values across turns and agents.

\subsection{Causal mediation analysis of gender bias}\label{lit:vig}
Vig, Gehrmann, Belinkov, Qian, Nevo, Singer and Shieber (2020) in ``Investigating gender bias in language models using causal mediation analysis''~\cite{vig2020causal} investigate which internal components of a language model are responsible for gender-biased predictions. The work adapts causal mediation analysis from the social sciences to neural networks, measuring the total effect of an input intervention (for example, changing a profession word) and decomposing it into direct and indirect effects mediated by individual neurons and attention heads. The study uses GPT-2 models of several sizes with templated professions prompts and the Winobias and Winogender coreference datasets. The authors evaluate the size and distribution of the mediated effects and report that gender bias effects are sparse, concentrated in a small number of components, notably attention heads in the middle layers, and that the total bias effect grows with model size. The main limitations are that the analysis covers single, static prompts and real gender terms, so learned and contextual bias cannot be separated, and that no intervention is proposed to reduce the bias. The work informs the present project through the idea that bias can be localised to identifiable heads, which the reservoir analysis tests on cached values during a conversation, and through the use of counterfactual interventions as evidence.

\subsection{Locating and editing factual associations}\label{lit:rome}
Meng, Bau, Andonian and Belinkov (2022) in ``Locating and editing factual associations in GPT''~\cite{meng2022rome} investigate where a language model stores factual associations such as the country in which a landmark lies. The work introduces causal tracing, which corrupts the input, restores the hidden state of one component at a time, and measures how much of the correct prediction is recovered. It finds that decisive computations take place in middle-layer feed-forward modules when the model processes the last token of the subject, and proposes Rank-One Model Editing (ROME), which changes a single fact by a rank-one update of one module's weights. The study uses GPT-2 XL and GPT-J and evaluates editing on the zsRE benchmark and on CounterFact, a dataset introduced in the paper. The authors assess efficacy, generalisation to paraphrases and specificity to unrelated facts, and report that ROME maintains both generalisation and specificity better than the fine-tuning and meta-learning editors it is compared with. The main limitations for the present problem are that ROME permanently modifies weights and that tracing is performed on single prompts. The work informs this project in two ways: its restore-and-measure logic is the model for the seed-matched replay used to estimate causal effects, and the contrast with weight editing defines the project's choice to edit only the runtime cache.

\subsection{Inference-time intervention for truthfulness}\label{lit:iti}
Li, Patel, Vi\'egas, Pfister and Wattenberg (2023) in ``Inference-time intervention: Eliciting truthful answers from a language model''~\cite{li2023iti} investigate whether a model's truthfulness can be improved at inference time without retraining. The work trains linear probes on the activations of every attention head to find heads whose outputs separate true from false answers, and then shifts the outputs of a small number of these heads along the truth-correlated directions during generation. The study uses LLaMA-family models, including the instruction-tuned Alpaca, and the TruthfulQA benchmark. The authors evaluate truthfulness and informativeness, and report that the intervention raises the truthfulness of Alpaca on TruthfulQA from 32.5\% to 65.1\%, while noting a trade-off between truthfulness and helpfulness that is controlled by the intervention strength. The main limitations are that the intervention acts on head outputs rather than on stored memory, that it is evaluated on single-turn question answering, and that stronger interventions reduce helpfulness. The work informs the present project closely: the reservoir analysis follows the same probe-then-intervene pattern but applies it to cached value vectors, and the reported trade-off motivates the coherence budget, which rejects any intervention whose bias reduction comes at the cost of general competence.

\subsection{Multi-agent debate}\label{lit:debate}
Du, Li, Torralba, Tenenbaum and Mordatch (2024) in ``Improving factuality and reasoning in language models through multiagent debate''~\cite{du2024debate} investigate whether several instances of a language model can improve their answers by discussing them. The work proposes a debate procedure in which multiple agents first answer independently and then, over several rounds, read each other's answers and revise their own until they converge on a common answer. The study uses black-box chat models accessed through an interface that exposes only text. The authors evaluate accuracy on arithmetic, grade-school mathematics (GSM8K), chess-move prediction, MMLU and the factual accuracy of generated biographies, and report that debate improves both reasoning and factual validity over single-agent baselines, with gains increasing with the number of agents and rounds. The main limitations are that the interaction is studied only at the level of exchanged text, that agents can converge on a shared wrong answer, and that no internal mechanism of influence is examined. The work informs the present project by establishing that agents strongly influence each other, which raises the question the project asks: whether influence, including the spread of bias, can also travel through shared memory below the level of text.

\subsection{Survey of bias and fairness in language models}\label{lit:survey}
Gallegos, Rossi, Barrow, Tanjim, Kim, Dernoncourt, Yu, Zhang and Ahmed (2024) in ``Bias and fairness in large language models: A survey''~\cite{gallegos2024bias} consolidate the definitions, metrics, datasets and mitigation techniques for social bias in language models. The work proposes three taxonomies: one of metrics, organised by whether they operate on embeddings, on model probabilities or on generated text; one of datasets, organised by whether they use counterfactual inputs or prompts; and one of mitigation techniques, organised by the stage at which they intervene, namely pre-processing, in-training, intra-processing at inference time, and post-processing. The study is a literature survey and does not run new experiments, so numerical findings are not applicable. The authors find that many metrics are only weakly correlated with each other and with downstream harms, and that most mitigation work concentrates on data and training. The main limitation for the present purpose is that the survey covers single-model, single-turn settings and does not consider bias that accumulates over a conversation or spreads between agents. The work informs the present project by placing its bias index among probability-based metrics and its cache interventions in the intra-processing category, which the survey identifies as comparatively under-explored.

\subsection{BBQ bias benchmark}\label{lit:bbq}
Parrish, Chen, Nangia, Padmakumar, Phang, Thompson, Htut and Bowman (2022) in ``BBQ: A hand-built bias benchmark for question answering''~\cite{parrish2022bbq} investigate how strongly question-answering models rely on social stereotypes. The work contributes a hand-built dataset of 58,492 examples organised in question sets that target attested biases against protected groups along nine social dimensions relevant to US English contexts. Each question is asked both with an ambiguous context, in which the correct answer is ``unknown'', and with a disambiguated context that states the answer. The study evaluates UnifiedQA, RoBERTa and DeBERTaV3 models fine-tuned for multiple-choice question answering. The authors measure accuracy and a bias score, and report that models rely strongly on stereotypes in ambiguous contexts and that, even with disambiguating information, accuracy is up to 3.4 percentage points higher when the correct answer matches a stereotype, rising to over 5 points for gender. The main limitations are the US-centred scope and the single-turn, static format, which cannot show how bias builds up in a conversation; with real group names, learned and contextual bias are also mixed. The work informs the present project through the ambiguous-context question format, which is reproduced in the project's probe bank and training corpus with fictional groups so that the direction of bias is known by construction.

\subsection{TruthfulQA}\label{lit:tqa}
Lin, Hilton and Evans (2022) in ``TruthfulQA: Measuring how models mimic human falsehoods''~\cite{lin2022truthfulqa} investigate whether language models repeat common human misconceptions. The work contributes a benchmark of 817 questions in 38 categories, including health, law, finance and politics, written so that some humans would answer them falsely because of a false belief. The study evaluates GPT-3, GPT-Neo/J, GPT-2 and T5-based models with human evaluation and an automated judge model fine-tuned for the task. The authors report that the best model was truthful on 58\% of questions against 94\% for humans, and that larger models were generally less truthful, an inverse-scaling result they attribute to imitation of human text. The main limitations are the single-turn format and the English, adversarially selected question set, which does not show whether false beliefs persist over a conversation. The work informs the present project through the factual probes: short questions with one correct answer, in the style of TruthfulQA, are asked after every turn to flag hallucinations and to measure factual accuracy, which forms part of the coherence budget used to accept or reject interventions.

\clearpage
\begin{landscape}
\section{Comparison of reviewed literature}
Table~\ref{tab:literature} compares the twelve reviewed works on the same attributes. Study context covers models, datasets and settings; evaluation criteria list the measures reported by the authors. ``Not applicable'' marks attributes that do not apply to a work, such as numerical findings for a survey.

\begingroup\small\setlength{\tabcolsep}{4pt}\renewcommand{\arraystretch}{1.25}
\begin{longtable}{|L{2.2cm}|L{2.6cm}|L{3.5cm}|L{2.7cm}|L{3.0cm}|L{2.9cm}|L{3.3cm}|}
\caption{Comparison of reviewed literature}\label{tab:literature}\\\hline
\textbf{Work and citation} & \textbf{Problem or objective} & \textbf{Approach and contribution} & \textbf{Study context} & \textbf{Evaluation criteria} & \textbf{Key findings} & \textbf{Limitations and gap}\\\hline\endfirsthead
\textbf{Work and citation} & \textbf{Problem or objective} & \textbf{Approach and contribution} & \textbf{Study context} & \textbf{Evaluation criteria} & \textbf{Key findings} & \textbf{Limitations and gap}\\\hline\endhead
Work 01~\cite{xiao2024sinks} Xiao et al., 2024 & Stable streaming generation with a bounded cache & StreamingLLM: keep attention-sink tokens plus a recent window; sink-token pre-training & Llama-2, MPT, Falcon, Pythia; long texts, streaming QA & Perplexity, accuracy, speed-up & Stable up to 4M tokens; up to 22.2$\times$ speed-up & Sinks used only for efficiency; semantic content and behavioural effect not studied\\\hline
Work 02~\cite{zhang2023h2o} Zhang et al., 2023 & Reduce KV cache memory during inference & H$_2$O: evict all but recent and heavy-hitter tokens; submodular formulation & OPT, LLaMA, GPT-NeoX; HELM and lm-eval tasks & Accuracy, ROUGE, throughput, latency & 20\% cache kept; up to 29$\times$ throughput gain & Importance only for eviction; no link to bias\\\hline
Work 03~\cite{liu2023scissorhands} Liu et al., 2023 & Fit the cache in a fixed budget & Persistence-of-importance hypothesis; keep pivotal tokens & OPT family; LM and few-shot tasks & Quality vs compression ratio & Up to 5$\times$ (20$\times$ with 4-bit) compression without quality loss & Persistence not tested across turns or agents\\\hline
Work 04~\cite{oren2024msrnn} Oren et al., 2024 & Relate transformers to RNNs & Cache as multi-state RNN state; TOVA eviction policy & LLaMA-2, Mistral; long-range tasks & Quality vs cache budget & Near full quality with one-eighth of the cache & Conceptual and efficiency focus; no semantic tracking\\\hline
Work 05~\cite{todd2024fv} Todd et al., 2024 & How in-context tasks are represented & Causal mediation finds heads carrying a function vector & GPT-J, Llama-2 up to 70B; many ICL tasks & Causal effect of inserting the vector & Few heads carry the task; vectors transfer zero-shot & Single forward pass; no evolving cache\\\hline
Work 06~\cite{vig2020causal} Vig et al., 2020 & Which components cause gender bias & Causal mediation analysis over neurons and heads & GPT-2 sizes; professions templates, Winobias, Winogender & Total, direct and indirect effects & Bias effects sparse, concentrated in mid-layer heads & Static prompts, real gender terms; no mitigation\\\hline
Work 07~\cite{meng2022rome} Meng et al., 2022 & Where facts are stored; how to edit them & Causal tracing; ROME rank-one weight edit & GPT-2 XL, GPT-J; zsRE, CounterFact & Efficacy, generalisation, specificity & Facts in mid-layer MLPs; edits generalise and stay specific & Weights modified permanently; single prompts\\\hline
Work 08~\cite{li2023iti} Li et al., 2023 & Improve truthfulness without training & Probe heads, then shift their outputs at inference & LLaMA, Alpaca; TruthfulQA & Truthfulness, informativeness & Alpaca truthfulness 32.5\% to 65.1\% & Acts on head outputs, single turn; helpfulness trade-off\\\hline
Work 09~\cite{du2024debate} Du et al., 2024 & Improve reasoning and factuality & Multi-round debate between model instances & Black-box chat models; arithmetic, GSM8K, chess, MMLU, biographies & Accuracy, factual validity & Debate beats single agents; gains grow with agents and rounds & Text-level only; no internal mechanism\\\hline
Work 10~\cite{gallegos2024bias} Gallegos et al., 2024 & Consolidate bias metrics and mitigation & Taxonomies of metrics, datasets and mitigation stages & Literature survey & Not applicable & Metrics weakly correlated; mitigation mostly pre- and in-training & Single-turn view; inference-time mitigation under-explored\\\hline
Work 11~\cite{parrish2022bbq} Parrish et al., 2022 & Measure stereotype reliance in QA & 58,492 hand-built ambiguous and disambiguated questions & UnifiedQA, RoBERTa, DeBERTaV3; nine social dimensions & Accuracy, bias score & Strong stereotype reliance; up to 3.4 points (over 5 for gender) accuracy gap & US-centred, single-turn; real groups mix learned and contextual bias\\\hline
Work 12~\cite{lin2022truthfulqa} Lin et al., 2022 & Measure imitation of human falsehoods & 817 adversarial questions in 38 categories & GPT-3, GPT-Neo/J, GPT-2, T5 & Truthfulness (human and model judge) & Best model 58\% vs human 94\%; larger models less truthful & Single-turn; persistence over dialogue not measured\\\hline
\end{longtable}\endgroup
\end{landscape}
\clearpage

\section{Summary of the research gap}
Read together, the reviewed works establish four facts that the project builds on and one gap that it addresses. The first four papers show that the KV cache has persistent, high-attention structure, but they use this structure only to save memory. The next four show that behaviour, including bias, can be localised to a few components and changed by targeted edits, but they analyse single prompts and edit weights or head outputs rather than stored memory. Multi-agent debate shows that agents change each other's answers, but only through text. The bias and truthfulness literature defines how to measure the behaviours of interest, but in static, single-turn settings. No reviewed work records the cache with turn, agent and channel provenance, separates a memory channel from a text channel when behaviour spreads between agents, or tests cache-only interventions against negative controls and a coherence budget. These are the capabilities developed in Chapter~5 and evaluated in Chapter~6.
